\renewcommand{\Anno}{2007}

\renewcommand{\NumeroProblema}{4}

\renewcommand{\Testo}{
La prossima edizione della ``Coppa America'' si svolgerà nelle acque davanti al litorale toscano.  La Scuola S.\ Anna ha deciso di parteciparvi con la sua barca ``Mascalzone Pisano''.  I tecnici hanno calcolato che la barca, quando procede lungo una direzione che forma un angolo $\theta$ con la direzione contraria a quella del vento, procede a velocità $v=v_{0}\sin^2(\theta/2)$, dove $v_{0}$ è la velocità del vento.  Ad esempio, se $\theta=0^{\circ}$, cioè se la barca prova ad andare esattamente controvento, allora $v=0$ (dunque è impossibile andare esattamente controvento); se $\theta=180^{\circ}$, cioè se la barca procede con il vento il poppa, allora $v=v_{0}$.
	
La parte fondamentale delle regate è costituita dai tratti di \emph{bolina}, che consistono nel percorrere una certa distanza $L$ in direzione contraria a quella del vento.
	
Il percorso della barca è una spezzata costituita da un certo numero di tratti rettilinei, ognuno dei quali viene percorso con la velocità costante data dalla formula calcolata dai tecnici. Inoltre ogni \emph{virata}, cioè ogni cambio di direzione della barca, fa perdere un certo tempo fisso $T_{0}$, indipendente dall'ampiezza del cambio di direzione.
	\begin{center}
		\psset{unit=5ex}
		\pspicture(-2,-1.7)(11,0.7)
		\psline(0,0)(10,0)
		\psdots(0,0)(10,0)
		\rput(5,-0.4){$L$}
		\rput(10,0.5){Fine bolina}
		\rput(0,0.5){Inizio bolina}
		\psline[linecolor=blue]{<-}(9,-1)(11,-1)
		\rput(10,-0.7){$v_{0}$}
		\rput(10,-1.5){Dir. vento}
		\psset{arrowsize=3pt 4.3}
		\psline[linecolor=green]{->}(0,0)(2,-1)
		\psline[linecolor=green]{->}(2,-1)(5,1)
		\psline[linecolor=green]{->}(5,1)(8,-1.5)
		\psline[linecolor=green]{->}(8,-1.5)(10,0)
		\endpspicture
	\end{center}
	
	\begin{enumerate}
		\item Nel caso in cui lo skipper decida di suddividere la bolina in 2 tratti di uguale lunghezza (dunque con una sola virata in mezzo), determinare il minimo tempo necessario per completare il percorso.
	
		\item Nel caso in cui lo skipper decida di suddividere la bolina in 2 tratti, eventualmente di lunghezza differente, determinare il minimo tempo necessario per completare il percorso.
	
		\item Determinare il minimo tempo per completare il tratto di bolina nel caso generale.
	\end{enumerate}
}

\renewcommand{\Risposta}{
In tutti e 3 i casi il tempo ottimale è $T_{0}+8L/v_{0}$, corrispondente a suddividere la bolina i 2 tratti di uguale lunghezza inclinati di $60^{\circ}$ rispetto alla direzione del vento.
}

\renewcommand{\Aiutino}{
Consideriamo la velocità orizzontale della barca, cioè la proiezione della sua velocità sulla congiungente Partenza/Arrivo. Mostrare che la velocità orizzontale è data da $v_{0}\cos\theta/\sin{\theta/2}$, dove $\theta$ indica l'angolo rispetto alla velocità del vento. Per quale valore di $\theta$ tale velocità è massima?
}

\renewcommand{\SoluzioneUfficiale}{\TitoloSoluzione \DomandaA 
Non potendo la barca andare completamente controvento, ed avendo lo skipper scelto di fare un'unica virata a metà percorso, la situazione è quella descritta nella seguente figura.
\begin{center}
			\psset{unit=5ex}
			\pspicture(-2,-2)(12,0.7)
			\psline(0,0)(10,0)
			\psdots(0,0)(10,0)(5,0)
			\rput[l](10.2,0){Fine bolina}
			\rput[r](-0.2,0){Inizio bolina}
			\rput(2.5,0.5){$L/2$}
			\rput(7.5,0.5){$L/2$}
			\psline[linecolor=blue]{<-}(9,-1)(11,-1)
			\rput(10,-0.7){$v_{0}$}
			\rput(10,-1.5){Dir. vento}
			\psset{arrowsize=3pt 4.3}
			\pnode(0,0){S}
			\pnode(10,0){B}
			\pnode(5,-2){P}
			\psline[linecolor=green]{->}(S)(P)
			\psline[linecolor=green]{->}(P)(B)
			\psline[linecolor=blue,linestyle=dashed](5,0)(P)
			\pstMarkAngle[LabelSep=1.5,MarkAngleRadius=0.8]{P}{S}{B}{$\theta$}
			\pstMarkAngle[LabelSep=1.5,MarkAngleRadius=0.8]{S}{B}{P}{$\theta$}
			\endpspicture
\end{center}

I 2 tratti, lunghi $L/(2\cos\theta)$, vengono percorsi a velocità $v_{0}\sin^{2}(\theta/2)$. Ne segue che il tempo di percorrenza è dato da
$$T=T_{0}+2\cdot\frac{L}{2\cos\theta}\cdot
\frac{1}{v_{0}\sin^{2}(\theta/2)}.$$

Minimizzare il tempo di percorrenza è pertanto equivalente a massimizzare la funzione
$$f(\theta)=\cos\theta\cdot\sin^{2}\frac{\theta}{2}=\frac{1}{2}\cos\theta
(1-\cos\theta)$$
al variare di $\theta$ tra 0 e $\pi/2$.  In tale intervallo $\cos\theta$ varia tra 0 e 1.  Ci si riconduce quindi a considerare la funzione $x(1-x)/2$ per $x$ tra 0 e $1$.  Tale funzione rappresenta una parabola rivolta verso il basso con massimo nel vertice, corrispondente ad $x=1/2$, dove vale $1/8$.  Risalendo, abbiamo che $f(\theta)$ assume il massimo quando $\cos\theta=1/2$, cioè (tenendo conto del vincolo $\theta\in[0,\pi/2]$) quando $\theta=\pi/3$.

Abbiamo pertanto dimostrato che 
\begin{equation}
    \cos\theta\cdot\sin^{2}(\theta/2)\leq\frac{1}{8}
    \quad\quad\forall \theta\in[0,\pi/2],
    \label{est-theta}
\end{equation}
con uguaglianza se e solo se $\theta=\pi/3$.  Ne segue che la strategia ottimale consiste nell'impostare i 2 tratti con un angolo di $60^{\circ}$ rispetto alla direzione del vento, completando la bolina in un tempo di $T_{0}+8L/v_{0}$.

\DomandaB Se la bolina viene suddivisa in 2 tratti eventualmente di lunghezza
diversa, la situazione è quella descritta nella seguente figura.
\begin{center}
			\psset{unit=5ex}
			\pspicture(-2,-2)(12,0.7)
			\psline(0,0)(10,0)
			\psdots(0,0)(10,0)(3,0)
			\rput[l](10.2,0){Fine bolina}
			\rput[r](-0.2,0){Inizio bolina}
			\rput(1.5,0.5){$\ell_{1}$}
			\rput(6.5,0.5){$\ell_{2}$}
			\psline[linecolor=blue]{<-}(9,-1)(11,-1)
			\rput(10,-0.7){$v_{0}$}
			\rput(10,-1.5){Dir. vento}
			\psset{arrowsize=3pt 4.3}
			\pnode(0,0){S}
			\pnode(10,0){B}
			\pnode(3,-2){P}
			\psline[linecolor=green]{->}(S)(P)
			\psline[linecolor=green]{->}(P)(B)
			\psline[linecolor=blue,linestyle=dashed](3,0)(P)
			\pstMarkAngle[LabelSep=1.2,MarkAngleRadius=0.8]{P}{S}{B}{$\theta_{1}$}
			\pstMarkAngle[LabelSep=2,MarkAngleRadius=1.2]{S}{B}{P}{$\theta_{2}$}
			\endpspicture
\end{center}

Gli stessi calcoli del caso precedente portano ora ad un tempo di percorrenza pari a
$$T=T_{0}+\frac{\ell_{1}}{\cos\theta_{1}}\cdot\frac{1}{v_{0}\sin^{2}\theta_{1}}+
\frac{\ell_{2}}{\cos\theta_{2}}\cdot\frac{1}{v_{0}\sin^{2}\theta_{2}}.$$

Ricordando la (\ref{est-theta}) possiamo minorare le 2 frazioni
ottenendo che
$$T\geq T_{0}+\frac{8\ell_{1}}{v_{0}}+\frac{8\ell_{2}}{v_{0}}\geq
T_{0}+\frac{8L}{v_{0}},$$
con uguaglianza se e solo se $\theta_{1}=\theta_{2}=60^{\circ}$, cioè se e solo se la rotta è quella descritta al punto (a).

Ad essere rigorosi, bisognerebbe anche escludere strategie esotiche in cui uno dei 2 tratti va all'indietro o perpendicolarmente rispetto alla direzione del vento, come nella seguente figura.
\begin{center}
		\psset{unit=5ex}
		\pspicture(-2,-2)(12,0.7)
		\pnode(0,0){S}
		\pnode(10,0){B}
		\pnode(-2,-2){P}
		\psline(0,0)(10,0)
		\psline[linecolor=blue](-2,0)(0,0)
		\psdots(0,0)(10,0)(-2,0)
		\rput[l](10.2,0){Fine bolina}
		%\rput[r](-0.2,0){Inizio bolina}
		\rput(-1,-0.3){$\ell_{1}$}
		\psline[linecolor=blue]{<-}(9,-1)(11,-1)
		\rput(10,-0.7){$v_{0}$}
		\rput(10,-1.5){Dir. vento}
		\pcline[offset=2ex,linewidth=0.5\pslinewidth,linecolor=cyan]{<->}(-2,0)(B)
		\lput*{U}{$\ell_{2}$}
		\psset{arrowsize=3pt 4.3}
		\psline[linecolor=green]{->}(S)(P)
		\psline[linecolor=green]{->}(P)(B)
		\psline[linecolor=blue,linestyle=dashed](-2,0)(P)	
		\pstMarkAngle[LabelSep=0.6,MarkAngleRadius=0.3]{P}{S}{B}{$\theta_{1}$}	
		\pstMarkAngle[LabelSep=3,MarkAngleRadius=1.8]{S}{B}{P}{$\theta_{2}$}
		\endpspicture
\end{center}

In tal caso comunque si ottiene la stima voluta considerando semplicemente i tratti in avanti.  Nel caso della figura, trascurando il tratto indietro, si ha comunque che
$$T>T_{0}+
\frac{\ell_{2}}{\cos\theta_{2}}\cdot\frac{1}{v_{0}\sin^{2}\theta_{2}}
\geq T_{0}+\frac{8\ell_{2}}{v_{0}}\geq T_{0}+\frac{8L}{v_{0}},$$
dunque senza possibilità di uguaglianza.

\DomandaC Consideriamo ora il caso in cui la bolina è suddivisa in
$k\geq 2$ tratti, che supponiamo per iniziare essere tutti ``in avanti''.
\begin{center}
			\psset{unit=5ex}
			\pspicture(-2,-3)(11,2)
			\psline(0,0)(10,0)
			\psdots(0,0)(10,0)(2,0)(4,0)(7,0)(8,0)
			\rput[l](10.2,0){Fine bolina}
			\rput[r](-0.2,0){Inizio bolina}
			\rput(1,-0.4){$\ell_{1}$}
			\rput(3,-0.4){$\ell_{2}$}
			\rput(5.5,0.4){$\ell_{3}$}
			\rput(7.5,0.4){$\ell_{4}$}
			\rput(9,0.4){$\ell_{5}$}
			\psline[linecolor=blue]{<-}(9,-1)(11,-1)
			\rput(10,-0.7){$v_{0}$}
			\rput(10,-1.5){Dir. vento}
			\psset{arrowsize=3pt 4.3}
			\pnode(0,0){S}
			\pnode(10,0){B}
			\pnode(3,-2){P}
			\psline[linecolor=green]{->}(0,0)(2,2)
			\psline[linecolor=green]{->}(2,2)(4,-1)
			\psline[linecolor=green]{->}(4,-1)(7,-3)
			\psline[linecolor=green]{->}(7,-3)(8,-1)
			\psline[linecolor=green]{->}(8,-1)(10,0)
			\psline[linecolor=blue,linestyle=dashed](2,0)(2,2)
			\psline[linecolor=blue,linestyle=dashed](4,0)(4,-1)
			\psline[linecolor=blue,linestyle=dashed](7,0)(7,-3)
			\psline[linecolor=blue,linestyle=dashed](8,0)(8,-1)
			\endpspicture
\end{center}

Ragionando come nel caso precedente avremo che 
\begin{eqnarray*}
    T & =  & (k-1)T_{0}+\sum_{i=1}^{k}
    \frac{\ell_{i}}{\cos\theta_{i}}\cdot\frac{1}{v_{0}\sin^{2}\theta_{i}}
    \\
     & \geq & (k-1)T_{0}+\frac{8}{v_{0}}\sum_{i=1}^{k}\ell_{i}  \\
     & \geq & (k-1)T_{0}+\frac{8L}{v_{0}}  \\
     & \geq & T_{0}+\frac{8L}{v_{0}}.
\end{eqnarray*}

Ancora una volta si ha uguaglianza se e solo se $k=2$ e $\theta_{1}=\theta_{2}=60^{\circ}$, cioè se e solo se siamo nella situazione descritta al punto (a).

Nel caso ci fossero tratti indietro o perpendicolari alla direzione del vento, basta ignorarli nella somma ed osservare che comunque la somma degli $\ell_{i}$ corrispondenti ai tratti in avanti è maggiore od uguale ad $L$.
}

\renewcommand{\SoluzioniAlternative}{\TitoloSoluzioneAlternativaA I
tre punti del problema potevano essere trattati in maniera unitaria ricorrendo al concetto di ``velocità orizzontale'', definita come la proiezione del vettore velocità sulla congiungente Partenza/Arrivo.

Quando la barca procede su una rotta che forma un angolo $\theta$ con la congiungente (si intende ora che $\theta\in[0,\pi]$, in cui angoli tra $\pi/2$ e $\pi$ rappresentano direzioni che tendono ad allontanare dall'arrivo) tale velocità orizzontale è un vettore di lunghezza $v_{0}\sin(\theta/2)\cos\theta$, dunque di lunghezza sempre minore od uguale di $v_{0}/8$ (e qui ancora una volta si utilizza che il massimo della funzione $f(\theta)$ è $1/8$).  È ora evidente che il tempo di percorrenza minimo si realizza quando tale vettore ha lunghezza $v_{0}/8$ durante tutto il tragitto, da cui si ricavano tutti i risultati precedenti in un colpo solo.

In quest'ottica è ancora più evidente l'assurdità delle strategie esotiche, le quali comprendono tratti in cui il vettore velocità è negativo.

}

\renewcommand{\Osservazioni}{\TitoloOsservazione Il massimo della
funzione $f(\theta)$ si poteva anche determinare ricorrendo al calcolo differenziale.  Calcolando la sua derivata prima si ha infatti che
$$f'(\theta)=-\sin\theta\sin^{2}\frac{\theta}{2}+
\cos\theta\cdot \sin\frac{\theta}{2}\cos\frac{\theta}{2}.$$

Ricordando le formule di duplicazione
$$\sin \theta=2\sin\frac{\theta}{2}\cos\frac{\theta}{2},
\hspace{5em}
\cos\theta=1-2\sin^{2}\frac{\theta}{2},$$
possiamo riscrivere $f'(\theta)$ nella forma
$$f'(\theta)=\sin\theta\left(\frac{1}{2}-
2\sin^{2}\frac{\theta}{2}\right).$$

Nell'intervallo $(0,\pi/2)$ tale derivata si annulla solo quando $\sin^{2}(\theta/2)=1/4$, e cioè per $\theta=\pi/3$.  Questo non basta per concludere che $\theta=\pi/3$ è il punto di massimo! Per esserne certi bisogna confrontare il valore in $\pi/3$ con i valori assunti agli estremi dell'intervallo.  Poiché $f$ vale 0 agli estremi dell'intervallo, e $f(\pi/3)=1/8$, possiamo concludere che $\pi/3$ è effettivamente il punto di massimo.

L'andamento di $f(\theta)$ è rappresentato nel seguente grafico.
\begin{center}
	\psset{xunit=2pt,yunit=10ex}
	\pspicture(-20,-1)(100,2)
	\psplot[linewidth=2\pslinewidth]{0}{90}{10 x cos x 2 div sin 2 exp mul mul}
	\psline[linewidth=0.5\pslinewidth]{->}(-10,0)(100,0)
	\psline[linewidth=0.5\pslinewidth]{->}(0,-0.5)(0,2)
	\rput[t](90,-0.1){$90^{\circ}$}
	\rput[t](60,-0.1){$60^{\circ}$}
	\rput[r](-5,1.25){$1/8$}
	\rput[tl](100,-0.1){$\theta$}
	\rput[l](4,2){$f(\theta)$}
	\psline[linecolor=magenta,linestyle=dashed](60,0)(60,1.25)
	\psline[linecolor=magenta,linestyle=dashed](0,1.25)(60,1.25)
	\endpspicture
\end{center}

Il calcolo risulta ancora più semplice se si scrive $f(\theta)$ nella forma $\cos\theta(1-\cos\theta)/2$.  In tal caso la derivata si scrive infatti come $f'(\theta)=\sin\theta(1-2\cos\theta)/2$ e si vede immediatamente che nell'intervallo $(0,\pi/2)$ l'unico punto in cui si annulla è quello in cui $\cos\theta=1/2$, da cui ancora una volta $\theta=\pi/3$.
}

\renewcommand{\Commenti}{
}